\hypertarget{appendix-2391db8052b1}{}
\section*{Cross products and oriented area}
\chapter{Cross products and oriented area}\label{appendix-2391db8052b1-cross-products-and-oriented-area}

\textbf{Curriculum depth:} Introduced · \textbf{First seen:} Phase6 · \textbf{Latest development:} Phase6

\section{Definition and mechanism}\label{appendix-2391db8052b1-definition-and-mechanism}

In oriented Euclidean three-dimensional space, the cross product is a perpendicular area vector. Its magnitude is parallelogram area; swapping inputs reverses its orientation. Ordinary component formulas assume an oriented orthonormal basis.

Panel edges \((4,0,0)\) and \((0,3,0)\) give \((0,0,12)\), so triangle area is \(6\). This is also the reconstructed Part 542 definition missing as a standalone supplied lesson.

\section{Prerequisites}\label{appendix-2391db8052b1-prerequisites}

\begin{itemize}
\tightlist
\item
  \hyperref[appendix-8bc4e8fc547c]{Dot products and orthogonality} --- Inner-product measurements supply the perpendicularity and alignment constraints.
\end{itemize}

\section{Used by}\label{appendix-2391db8052b1-used-by}

\begin{itemize}
\item
  \href{../01-Foundations-I/main.pdf\#nameddest=appendix-da478be01d02}{Determinants orientation and volume (01-Foundations-I, appendix-da478be01d02)}
\item
  \href{../01-Foundations-I/main.pdf\#nameddest=appendix-965d9260eeab}{Quaternion algebra and Hamilton product (01-Foundations-I, appendix-965d9260eeab)}
\item
  \hyperref[appendix-87d7626c8d37]{Coordinate handedness and orientation} --- Builds on this mechanism in the connected geometry story.
\item
  \hyperref[appendix-5caa058bd9c7]{Triple products and signed volume} --- Builds on this mechanism in the connected geometry story.
\item
  \hyperref[appendix-e569e60e18e6]{Robust vector angles} --- Builds on this mechanism in the connected geometry story.
\item
  \hyperref[appendix-e31893dec985]{Surface tangents normals and bitangents} --- Builds on this mechanism in the connected geometry story.
\item
  \hyperref[appendix-341296fb3f64]{Robust geometric predicates} --- Builds on this mechanism in the connected geometry story.
\item
  \hyperref[appendix-e1960925d1b8]{Exterior products and geometric algebra} --- Builds on this mechanism in the connected geometry story.
\item
  \href{../03-Visual-Computing-I/main.pdf\#nameddest=appendix-7107009c505b}{Local rigidity and mesh orientation fields (03-Visual-Computing-I, appendix-7107009c505b)} \#\# Recall and next depth
\end{itemize}

Prove perpendicularity and the area identity; distinguish a two-dimensional signed area from a vector cross. Full production proofs, target-specific integration and learner mastery require further evidence.

\section{Source and navigation}\label{appendix-2391db8052b1-source-and-navigation}

\href{Continuous\%20Notes/Computing\%20Foundations\%20-\%20From\%20Bits\%20to\%20C++\%20and\%20Systems\%20APIs\#Oriented\%20area\%20tells\%20us\%20which\%20side\%20of\%20the\%20panel\%20is\%20front}{Teaching narrative} · \href{../01-Foundations-I/main.pdf\#nameddest=appendix-fd090d982a15-phase-6---parts-501-600}{Coverage and source gap (01-Foundations-I, appendix-fd090d982a15-phase-6-\/-\/-parts-501-600)} · \href{../01-Foundations-I/main.pdf\#nameddest=appendix-a103e7846d21-phase-6---source-analysis-and-verification}{Evidence and limits (01-Foundations-I, appendix-a103e7846d21-phase-6-\/-\/-source-analysis-and-verification)} · \hyperref[appendix-4ad47bdd269a]{Dependency Map}
